\pagestyle{empty}
\tight
\begin{tblr}{width=\textwidth,colspec={|Q[c,m,1.9cm]|X[l,m]|},hlines,vlines,hspan=minimal,row{1}={bg=headblue},rowsep=1pt,colsep=3pt}
\SetCell{c,m}\hfil\logo\hfil & \textbf{POLITEKNIK NEGERI MALANG}\par\textbf{JURUSAN TEKNOLOGI INFORMASI}\par\textbf{PROGRAM STUDI: D4 TEKNIK INFORMATIKA} \\
\SetCell[c=2]{c} \textbf{RENCANA PEMBELAJARAN SEMESTER (RPS)} & \\
\end{tblr}
\begin{longtblr}{width=\textwidth,colspec={|Q[l,m,1.9cm]|X[0.85,l,m]|X[1.45,l,m]|X[0.75,l,m]|X[1.15,l,m]|},hlines,vlines,hspan=minimal,rowsep=1pt,colsep=2pt,row{1,3}={bg=headgray,font=\bfseries}}
MATA KULIAH & KODE & BOBOT (SKS)/jam & SEMESTER & TGL. PENYUSUNAN \\
Workshop Teknologi Terapan & RTI256205 & 3 SKS/6 jam & 6 & 1 Juli 2025 \\
OTORISASI & Dosen Pengembang RPS & Koordinator RMK & \SetCell[c=2]{l} Ka PRODI &  \\
 & Triana Fatmawati, S.T., M.T. & Triana Fatmawati, S.T., M.T. & \SetCell[c=2]{l}{\vspace{8mm}\par Dr. Ely Setyo Astuti, S.T., M.T.} &  \\
\SetCell[r=11]{l} Capaian Pembelajaran (CP) & \SetCell[c=4]{l,bg=headgray,font=\bfseries} Capaian Pembelajaran Lulusan Program Studi (CPL-Prodi) &  &  &  \\
 & CPL06 & \SetCell[c=3]{l} Mampu merancang, mengimplementasikan, menguji, melakukan deployment, memelihara, serta menjamin mutu perangkat lunak yang menjawab permasalahan dan memenuhi kebutuhan pengguna. &  &  \\
 & CPL08 & \SetCell[c=3]{l} Mampu mengelola proyek teknologi informasi secara profesional melalui perencanaan, koordinasi, komunikasi, dan optimalisasi sumber daya. &  &  \\
 & \SetCell[c=4]{l,bg=headgray,font=\bfseries} Capaian Pembelajaran mata kuliah (CPL-MK) &  &  &  \\
 & CPMK0603 & \SetCell[c=3]{l} Mahasiswa mampu menerapkan teknik pengujian perangkat lunak dan quality assurance dalam konteks industri nyata. &  &  \\
 & CPMK0802 & \SetCell[c=3]{l} Mahasiswa mampu mengimplementasikan siklus hidup rekayasa perangkat lunak secara menyeluruh dalam workshop industri. &  &  \\
 & CPMK0803 & \SetCell[c=3]{l} Mahasiswa mampu mengeksekusi dan mengelola proyek workshop teknologi secara profesional dalam tim lintas fungsi. &  &  \\
 & \SetCell[c=4]{l,bg=headgray,font=\bfseries} Kemampuan akhir tiap tahapan belajar (Sub-CPMK) &  &  &  \\
 & SCPMK0603-05401 & \SetCell[c=3]{l} Mahasiswa mampu merancang dan mengeksekusi rencana pengujian komprehensif termasuk unit testing, integration testing, dan UAT. &  &  \\
 & SCPMK0802-05402 & \SetCell[c=3]{l} Mahasiswa mampu mengembangkan produk teknologi terapan melalui iterasi sprint berbasis kebutuhan industri nyata. &  &  \\
 & SCPMK0803-05403 & \SetCell[c=3]{l} Mahasiswa mampu memimpin tim dalam workshop, mengelola backlog, dan melaporkan kemajuan kepada stakeholder industri. &  &  \\
\SetCell[r=2]{l} Penilaian & \SetCell[c=4]{l} *Nilai CPMK didapat dari agregasi nilai SUB-CPMK yang dibebankan &  &  &  \\
 & \SetCell[c=4]{l}{\tiny\begin{tblr}{width=\linewidth,colspec={|Q[c,m,0.1\linewidth]|Q[c,m,0.16\linewidth]|X[l,m]|X[l,m]|X[l,m]|X[l,m]|X[l,m]|X[l,m]|Q[c,m,0.08\linewidth]|},hlines,vlines,hspan=minimal,rowsep=1pt,colsep=0.7pt,row{1,2}={bg=headgray,font=\bfseries}}
Kode CPMK & Kode SCPMK & \SetCell[c=6]{c} Bobot per Bentuk Penilaian & & & & & & Total Bobot \\
 & & Rencana & Sprint 1-2 & UTS & Sprint 3-4 & Laporan QA & UAS Demo & \\
CPMK0603 & SCPMK0603-05401 & 0\% & 8\%: rencana uji komprehensif & 0\% & 10\%: pengujian sprint 3-4 & 10\%: laporan QA final & 0\% & 28\% \\
CPMK0802 & SCPMK0802-05402 & 7\%: rencana pengembangan produk & 12\%: implementasi sprint 1-2 & 0\% & 10\%: implementasi sprint 3-4 & 0\% & 10\%: demo produk & 39\% \\
CPMK0803 & SCPMK0803-05403 & 8\%: rencana dan WBS workshop & 0\% & 15\%: milestone review & 0\% & 0\% & 10\%: defense kepada industri & 33\% \\
\SetCell[c=8]{r}\textbf{Total Per Penilaian} & & & & & & & & \textbf{100\%} \\
\end{tblr}} &  &  &  \\
Diskripsi Singkat Mata Kuliah & \SetCell[c=4]{l} Workshop Teknologi Terapan adalah capstone jalur magang/industri yang mengintegrasikan rekayasa perangkat lunak, quality assurance, dan manajemen proyek dalam mengembangkan produk teknologi nyata sesuai kebutuhan industri melalui pendekatan workshop berbasis sprint. &  &  &  \\
Bahan Kajian / Materi Pembelajaran & \SetCell[c=4]{l}\begin{enumerate}\item Rekayasa perangkat lunak dalam konteks industri: sprint planning dan backlog management\item Pengembangan produk iteratif: implementasi fitur, integrasi, dan validasi\item Quality assurance: unit testing, integration testing, dan UAT\item Manajemen proyek workshop: stakeholder communication, demo, dan retrospective\end{enumerate} &  &  &  \\
\SetCell[r=2]{l} Pustaka & \SetCell[c=4]{l} \textbf{Utama:}\begin{enumerate}\item Humble, J. 2011. Continuous Delivery. Pearson.\item Cohn, M. 2009. Succeeding with Agile. Addison-Wesley.\item Myers, G. 2011. The Art of Software Testing. Wiley.\end{enumerate} &  &  &  \\
 & \SetCell[c=4]{l} \textbf{Pendukung:} Materi LMS, dokumentasi resmi perangkat lunak, dan jobsheet praktikum. &  &  &  \\
Media Pembelajaran & \SetCell[c=2]{l} \textbf{Software:} IDE, Git, CI/CD tools, project management tools, LMS. &  & \SetCell[c=2]{l} \textbf{Hardware:} Komputer/laptop, proyektor. &  \\
Nama Dosen Pengampu & \SetCell[c=4]{l} Tim Pengajar Program Studi D4 TI &  &  &  \\
Matakuliah Syarat & \SetCell[c=4]{l} - &  &  &  \\
\end{longtblr}
\begin{longtblr}{width=\textwidth,colspec={|Q[c,m,0.06\textwidth]|X[1.35,l,m]|X[1.08,l,m]|X[1.16,l,m]|X[0.7,l,m]|X[1.24,l,m]|X[1.06,l,m]|X[0.98,l,m]|Q[c,m,0.05\textwidth]|},hlines,vlines,hspan=minimal,rowhead=2,row{1,2}={bg=headgreen,font=\bfseries},rowsep=1pt,colsep=1.5pt}
Mgg.\par Ke & Kemampuan Akhir Yang Direncanakan (Sub-CP-MK) & Bahan kajian (Materi Pembelajaran) & Bentuk dan Metode Pembelajaran & Estimasi Waktu & Pengalaman Belajar Mahasiswa & Kriteria \& Bentuk Penilaian & Indikator Penilaian & Bobot Penilaian (\%) \\
(1) & (2) & (3) & (4) & (5) & (6) & (7) & (8) & (9) \\
1 & SCPMK0803-05403 & Orientasi workshop industri dan pembentukan tim. & Modalitas: Blended Learning. Bentuk: Luring/Daring. Metode: Workshop berbasis sprint, peer collaboration, stakeholder review, dan retrospective. & 1 x 6 x 50' workshop industri; 1 x 6 x 50' dokumentasi dan tugas mandiri. & Mahasiswa mengerjakan aktivitas workshop untuk orientasi dan pembentukan tim serta menunjukkan evidence hasil belajar. & Bentuk penilaian: Orientasi Workshop Industri. Mengacu rubrik RTM. & Ketepatan implementasi; kualitas produk; validasi dan dokumentasi. & 3 \\
2 & SCPMK0802-05402 & Analisis kebutuhan industri dan requirement engineering. & Modalitas: Blended Learning. Bentuk: Luring/Daring. Metode: Workshop berbasis sprint, peer collaboration, stakeholder review, dan retrospective. & 1 x 6 x 50' workshop industri; 1 x 6 x 50' dokumentasi dan tugas mandiri. & Mahasiswa mengerjakan aktivitas workshop untuk analisis kebutuhan industri dan menunjukkan evidence hasil belajar. & Bentuk penilaian: Analisis Kebutuhan Industri. Mengacu rubrik RTM. & Ketepatan implementasi; kualitas produk; validasi dan dokumentasi. & 3 \\
3 & SCPMK0803-05403 & Perencanaan workshop: WBS, timeline, target deliverable. & Modalitas: Blended Learning. Bentuk: Luring/Daring. Metode: Workshop berbasis sprint, peer collaboration, stakeholder review, dan retrospective. & 1 x 6 x 50' workshop industri; 1 x 6 x 50' dokumentasi dan tugas mandiri. & Mahasiswa mengerjakan aktivitas workshop untuk perencanaan WBS dan timeline serta menunjukkan evidence hasil belajar. & Bentuk penilaian: Perencanaan Workshop. Mengacu rubrik RTM. & Ketepatan implementasi; kualitas produk; validasi dan dokumentasi. & 4 \\
4 & SCPMK0802-05402 & Setup environment dan architecture design. & Modalitas: Blended Learning. Bentuk: Luring/Daring. Metode: Workshop berbasis sprint, peer collaboration, stakeholder review, dan retrospective. & 1 x 6 x 50' workshop industri; 1 x 6 x 50' dokumentasi dan tugas mandiri. & Mahasiswa mengerjakan aktivitas workshop untuk setup environment dan architecture design serta menunjukkan evidence hasil belajar. & Bentuk penilaian: Setup Environment dan Architecture Design. Mengacu rubrik RTM. & Ketepatan implementasi; kualitas produk; validasi dan dokumentasi. & 4 \\
5 & SCPMK0802-05402 & Sprint 1: implementasi fitur core. & Modalitas: Blended Learning. Bentuk: Luring/Daring. Metode: Workshop berbasis sprint, peer collaboration, stakeholder review, dan retrospective. & 1 x 6 x 50' workshop industri; 1 x 6 x 50' dokumentasi dan tugas mandiri. & Mahasiswa mengerjakan aktivitas workshop untuk implementasi fitur core sprint 1 dan menunjukkan evidence hasil belajar. & Bentuk penilaian: Sprint 1 Implementasi Fitur Core. Mengacu rubrik RTM. & Ketepatan implementasi; kualitas produk; validasi dan dokumentasi. & 5 \\
6 & SCPMK0603-05401 & Sprint 1: testing dan dokumentasi. & Modalitas: Blended Learning. Bentuk: Luring/Daring. Metode: Workshop berbasis sprint, peer collaboration, stakeholder review, dan retrospective. & 1 x 6 x 50' workshop industri; 1 x 6 x 50' dokumentasi dan tugas mandiri. & Mahasiswa mengerjakan aktivitas workshop untuk testing dan dokumentasi sprint 1 dan menunjukkan evidence hasil belajar. & Bentuk penilaian: Sprint 1 Testing dan Dokumentasi. Mengacu rubrik RTM. & Ketepatan implementasi; kualitas produk; validasi dan dokumentasi. & 5 \\
7 & SCPMK0803-05403 & Sprint 1 review, demo, dan retrospective. & Modalitas: Blended Learning. Bentuk: Luring/Daring. Metode: Workshop berbasis sprint, peer collaboration, stakeholder review, dan retrospective. & 1 x 6 x 50' workshop industri; 1 x 6 x 50' dokumentasi dan tugas mandiri. & Mahasiswa mengerjakan aktivitas workshop untuk sprint review, demo, dan retrospective serta menunjukkan evidence hasil belajar. & Bentuk penilaian: Sprint 1 Review Demo Retrospective. Mengacu rubrik RTM. & Ketepatan implementasi; kualitas produk; validasi dan dokumentasi. & 5 \\
8 & SCPMK0803-05403 & UTS Milestone: evaluasi progres dan stakeholder review. & Modalitas: Blended Learning. Bentuk: Luring/Daring. Metode: Workshop berbasis sprint, peer collaboration, stakeholder review, dan retrospective. & 1 x 6 x 50' workshop industri; 1 x 6 x 50' dokumentasi dan tugas mandiri. & Mahasiswa mengerjakan aktivitas evaluasi progres dan stakeholder review serta menunjukkan evidence hasil belajar. & Bentuk penilaian: UTS Milestone Workshop. Mengacu rubrik RTM. & Ketepatan implementasi; kualitas produk; validasi dan dokumentasi. & 15 \\
9 & SCPMK0802-05402 & Sprint 2: pengembangan fitur lanjutan. & Modalitas: Blended Learning. Bentuk: Luring/Daring. Metode: Workshop berbasis sprint, peer collaboration, stakeholder review, dan retrospective. & 1 x 6 x 50' workshop industri; 1 x 6 x 50' dokumentasi dan tugas mandiri. & Mahasiswa mengerjakan aktivitas workshop untuk pengembangan fitur lanjutan sprint 2 dan menunjukkan evidence hasil belajar. & Bentuk penilaian: Sprint 2 Pengembangan Fitur Lanjutan. Mengacu rubrik RTM. & Ketepatan implementasi; kualitas produk; validasi dan dokumentasi. & 4 \\
10 & SCPMK0603-05401 & Sprint 2: integrasi dan regression testing. & Modalitas: Blended Learning. Bentuk: Luring/Daring. Metode: Workshop berbasis sprint, peer collaboration, stakeholder review, dan retrospective. & 1 x 6 x 50' workshop industri; 1 x 6 x 50' dokumentasi dan tugas mandiri. & Mahasiswa mengerjakan aktivitas workshop untuk integrasi dan regression testing sprint 2 serta menunjukkan evidence hasil belajar. & Bentuk penilaian: Sprint 2 Integrasi dan Regression Testing. Mengacu rubrik RTM. & Ketepatan implementasi; kualitas produk; validasi dan dokumentasi. & 5 \\
11 & SCPMK0803-05403 & Sprint 2 review dan Sprint 3 planning. & Modalitas: Blended Learning. Bentuk: Luring/Daring. Metode: Workshop berbasis sprint, peer collaboration, stakeholder review, dan retrospective. & 1 x 6 x 50' workshop industri; 1 x 6 x 50' dokumentasi dan tugas mandiri. & Mahasiswa mengerjakan aktivitas workshop untuk sprint 2 review dan sprint 3 planning serta menunjukkan evidence hasil belajar. & Bentuk penilaian: Sprint 2 Review Sprint 3 Planning. Mengacu rubrik RTM. & Ketepatan implementasi; kualitas produk; validasi dan dokumentasi. & 4 \\
12 & SCPMK0603-05401 & Sprint 3: finalisasi fitur dan performance testing. & Modalitas: Blended Learning. Bentuk: Luring/Daring. Metode: Workshop berbasis sprint, peer collaboration, stakeholder review, dan retrospective. & 1 x 6 x 50' workshop industri; 1 x 6 x 50' dokumentasi dan tugas mandiri. & Mahasiswa mengerjakan aktivitas workshop untuk finalisasi fitur dan performance testing sprint 3 serta menunjukkan evidence hasil belajar. & Bentuk penilaian: Sprint 3 Finalisasi dan Performance Testing. Mengacu rubrik RTM. & Ketepatan implementasi; kualitas produk; validasi dan dokumentasi. & 5 \\
13 & SCPMK0603-05401 & Sprint 3: UAT dan feedback incorporation. & Modalitas: Blended Learning. Bentuk: Luring/Daring. Metode: Workshop berbasis sprint, peer collaboration, stakeholder review, dan retrospective. & 1 x 6 x 50' workshop industri; 1 x 6 x 50' dokumentasi dan tugas mandiri. & Mahasiswa mengerjakan aktivitas workshop untuk UAT dan feedback incorporation sprint 3 serta menunjukkan evidence hasil belajar. & Bentuk penilaian: Sprint 3 UAT dan Feedback. Mengacu rubrik RTM. & Ketepatan implementasi; kualitas produk; validasi dan dokumentasi. & 4 \\
14 & SCPMK0603-05401 & Laporan QA dan dokumentasi testing final. & Modalitas: Blended Learning. Bentuk: Luring/Daring. Metode: Workshop berbasis sprint, peer collaboration, stakeholder review, dan retrospective. & 1 x 6 x 50' workshop industri; 1 x 6 x 50' dokumentasi dan tugas mandiri. & Mahasiswa mengerjakan aktivitas workshop untuk laporan QA dan dokumentasi testing final serta menunjukkan evidence hasil belajar. & Bentuk penilaian: Laporan QA dan Dokumentasi Testing. Mengacu rubrik RTM. & Ketepatan implementasi; kualitas produk; validasi dan dokumentasi. & 5 \\
15 & SCPMK0802-05402 & Sprint 4: product polish dan deployment. & Modalitas: Blended Learning. Bentuk: Luring/Daring. Metode: Workshop berbasis sprint, peer collaboration, stakeholder review, dan retrospective. & 1 x 6 x 50' workshop industri; 1 x 6 x 50' dokumentasi dan tugas mandiri. & Mahasiswa mengerjakan aktivitas workshop untuk product polish dan deployment sprint 4 serta menunjukkan evidence hasil belajar. & Bentuk penilaian: Sprint 4 Product Polish dan Deployment. Mengacu rubrik RTM. & Ketepatan implementasi; kualitas produk; validasi dan dokumentasi. & 5 \\
16 & SCPMK0803-05403 & Persiapan demo dan presentasi industry stakeholder. & Modalitas: Blended Learning. Bentuk: Luring/Daring. Metode: Workshop berbasis sprint, peer collaboration, stakeholder review, dan retrospective. & 1 x 6 x 50' workshop industri; 1 x 6 x 50' dokumentasi dan tugas mandiri. & Mahasiswa mengerjakan aktivitas workshop untuk persiapan demo dan presentasi industry stakeholder serta menunjukkan evidence hasil belajar. & Bentuk penilaian: Persiapan Demo dan Presentasi. Mengacu rubrik RTM. & Ketepatan implementasi; kualitas produk; validasi dan dokumentasi. & 4 \\
17 & SCPMK0803-05403, SCPMK0802-05402 & UAS: demo produk dan defense kepada industri. & Modalitas: Blended Learning. Bentuk: Luring/Daring. Metode: Workshop berbasis sprint, peer collaboration, stakeholder review, dan retrospective. & 1 x 6 x 50' workshop industri; 1 x 6 x 50' dokumentasi dan tugas mandiri. & Mahasiswa mengerjakan aktivitas UAS berupa demo produk dan defense kepada industri serta menunjukkan evidence hasil belajar. & Bentuk penilaian: UAS Demo Produk dan Defense. Mengacu rubrik RTM. & Ketepatan implementasi; kualitas produk; validasi dan dokumentasi. & 20 \\
\end{longtblr}
